\section{Implications for ReDoS}
\label{sec:redos}

ReDoS arises when a regular-expression engine spends superlinear or exponential
time exploring alternatives.  The bridge developed here suggests a semantic
pipeline for certified ReDoS analysis.

\paragraph{Step 1: Compile to a position automaton.}
For the regular core, this is syntax-directed and already implemented.  For
larger JavaScript fragments, compilation must account for captures and
zero-width assertions, or else restrict the analysis to a conservative regular
approximation.

\paragraph{Step 2: Classify ambiguity.}
Weber--Seidl criteria can identify unambiguous, finitely ambiguous,
polynomially ambiguous, and exponentially ambiguous automata.  In a verifier,
the analyzer can emit a graph certificate and Coq can check that the certificate
implies the claimed ambiguity class.

\paragraph{Step 3: Interpret the class as a trace bound.}
The trace/run theorem turns ambiguity bounds into bounds on successful
alternatives.  For many practical patterns, these alternatives are the source of
the engine's branching behavior.

\paragraph{Step 4: Account for failure.}
A full ReDoS detector must also bound failing branches.  We propose to model
failing branches through \emph{prefix ambiguity}: the number of partial runs that
consume a prefix and reach a live but ultimately rejecting configuration.  This
quantity is naturally expressible as ambiguity in product automata that track
both progress and residual obligations.  Formalizing this prefix layer is the
main technical extension beyond the successful-leaf correspondence.
